% ============================================================
% 测试卷 — 第6章 平面向量及其应用
% ============================================================
\documentclass[12pt, a4paper]{ctexart}
\usepackage[margin=2cm]{geometry}
\usepackage{amsmath, amssymb}
\usepackage{enumitem}
\usepackage{fancyhdr}
\usepackage{xcolor}
\usepackage{tikz}
\usepackage{graphicx}
\usepackage{tabularx}
\usepackage{multicol}

\definecolor{mainblue}{RGB}{41, 98, 168}

\pagestyle{fancy}
\fancyhf{}
\fancyhead[L]{\small\textcolor{mainblue}{2025-2026学年度·第6章测试卷}}
\fancyhead[R]{\small\textcolor{mainblue}{数学（必修第二册）}}
\fancyfoot[C]{\small —\ \thepage\ —}
\renewcommand{\headrulewidth}{0.6pt}
\renewcommand{\headrule}{\hbox to\headwidth{\color{mainblue}\leaders\hrule height \headrulewidth\hfill}}

\setCJKmainfont{Hiragino Sans GB}
\setCJKsansfont{Heiti SC}

\newcommand{\blank}[1]{\underline{\hspace{#1}}}

\begin{document}

\begin{center}
{\Large\bfseries 第6章\quad 平面向量及其应用\quad 测试卷}\\[6pt]
{\normalsize 数学（必修第二册）}\\[4pt]
{\small 考试时间：120分钟\quad 满分：150分}
\end{center}

\vspace{2pt}
\begin{center}
\begin{tabular}{|c|c|c|c|c|}
\hline
题号 & 一 & 二 & 三 & 四 & 总分 \\ \hline
得分 & & & & & \\ \hline
\end{tabular}
\end{center}
\vspace{4pt}

\noindent\textbf{注意事项：}
\begin{enumerate}[leftmargin=2em]
  \item 答题前，请在答题卡上填写姓名、班级和考号。
  \item 选择题请用2B铅笔在答题卡上涂黑对应选项。
  \item 非选择题请用黑色签字笔在答题卡指定区域作答。
\end{enumerate}

\vspace{6pt}

% ============ 一、单项选择题 ============
\noindent\rule{\textwidth}{0.4pt}\\[2pt]
{\large\bfseries 一、单项选择题}（本大题共8小题，每小题5分，共40分。每小题只有一个选项符合题目要求。）

\vspace{6pt}

\begin{enumerate}
\item 下列关于零向量的说法正确的是（\quad）
  \begin{enumerate}[label=\Alph*., itemsep=0pt]
    \item 零向量没有方向 \quad \item 零向量的模为 $0$
    \item 零向量与任一向量都不平行 \quad \item 零向量与任一向量都不相等
  \end{enumerate}

\item 设 $\vec{a}=(2,-3)$，$\vec{b}=(-4,6)$，则 $\vec{a}$ 与 $\vec{b}$ 的关系是（\quad）
  \begin{enumerate}[label=\Alph*., itemsep=0pt]
    \item 相等 \quad \item 共线且同向 \quad \item 共线且反向 \quad \item 不共线
  \end{enumerate}

\item 在 $\triangle ABC$ 中，$\overrightarrow{AB}=\vec{a}$，$\overrightarrow{AC}=\vec{b}$，则 $\overrightarrow{BC}=$（\quad）
  \begin{enumerate}[label=\Alph*., itemsep=0pt]
    \item $\vec{a}+\vec{b}$ \quad \item $\vec{a}-\vec{b}$ \quad \item $\vec{b}-\vec{a}$ \quad \item $-\vec{a}-\vec{b}$
  \end{enumerate}

\item 已知 $\vec{a}=(1,2)$，$\vec{b}=(3,-1)$，则 $2\vec{a}+\vec{b}=$（\quad）
  \begin{enumerate}[label=\Alph*., itemsep=0pt]
    \item $(5,3)$ \quad \item $(7,0)$ \quad \item $(5,0)$ \quad \item $(7,3)$
  \end{enumerate}

\item 已知 $\vec{a}=(1,2)$，$\vec{b}=(2,y)$，若 $\vec{a}\parallel\vec{b}$，则 $y=$（\quad）
  \begin{enumerate}[label=\Alph*., itemsep=0pt]
    \item $1$ \quad \item $2$ \quad \item $3$ \quad \item $4$
  \end{enumerate}

\item 已知 $\vec{a}=(1,2)$，$\vec{b}=(m,1)$，若 $\vec{a}\perp\vec{b}$，则 $m=$（\quad）
  \begin{enumerate}[label=\Alph*., itemsep=0pt]
    \item $-2$ \quad \item $2$ \quad \item $\dfrac{1}{2}$ \quad \item $-\dfrac{1}{2}$
  \end{enumerate}

\item 已知 $|\vec{a}|=3$，$|\vec{b}|=4$，$\vec{a}\cdot\vec{b}=-6$，则 $\vec{a}$ 与 $\vec{b}$ 的夹角为（\quad）
  \begin{enumerate}[label=\Alph*., itemsep=0pt]
    \item $\dfrac{\pi}{6}$ \quad \item $\dfrac{\pi}{3}$ \quad \item $\dfrac{2\pi}{3}$ \quad \item $\dfrac{5\pi}{6}$
  \end{enumerate}

\item 已知向量 $\vec{a}=(1,\sqrt{3})$，$\vec{b}=(0,2)$，则 $\vec{a}$ 在 $\vec{b}$ 方向上的投影为（\quad）
  \begin{enumerate}[label=\Alph*., itemsep=0pt]
    \item $1$ \quad \item $\sqrt{3}$ \quad \item $\dfrac{\sqrt{3}}{2}$ \quad \item $2$
  \end{enumerate}
\end{enumerate}

\vspace{6pt}

% ============ 二、多项选择题 ============
\noindent\rule{\textwidth}{0.4pt}\\[2pt]
{\large\bfseries 二、多项选择题}（本大题共4小题，每小题5分，共20分。全部选对得5分，部分选对得2分，有选错的得0分。）

\vspace{6pt}

\begin{enumerate}[resume]
\item 设 $\vec{a}$，$\vec{b}$ 为非零向量，下列命题正确的有（\quad）
  \begin{enumerate}[label=\Alph*., itemsep=0pt]
    \item 若 $\vec{a}\parallel\vec{b}$，则存在唯一实数 $\lambda$ 使得 $\vec{b}=\lambda\vec{a}$
    \item 若 $\vec{a}\parallel\vec{b}$ 且 $\vec{b}\parallel\vec{c}$，则 $\vec{a}\parallel\vec{c}$
    \item 若 $|\vec{a}+\vec{b}|=|\vec{a}|+|\vec{b}|$，则 $\vec{a}$ 与 $\vec{b}$ 同向
    \item $|\vec{a}|-|\vec{b}|\leqslant|\vec{a}+\vec{b}|\leqslant|\vec{a}|+|\vec{b}|$
  \end{enumerate}

\item 已知 $\vec{a}=(1,1)$，$\vec{b}=(-1,2)$，则下列结论正确的是（\quad）
  \begin{enumerate}[label=\Alph*., itemsep=0pt]
    \item $\vec{a}\cdot\vec{b}=1$
    \item $\vec{a}\perp\vec{b}$
    \item $|\vec{a}|=\sqrt{2}$
    \item $\vec{a}$ 与 $\vec{b}$ 的夹角为 $\dfrac{\pi}{4}$
  \end{enumerate}

\item 在正方形 $ABCD$ 中，设 $\overrightarrow{AB}=\vec{a}$，$\overrightarrow{AD}=\vec{b}$，$E$ 为 $BC$ 的中点，$F$ 为 $CD$ 的中点，则（\quad）
  \begin{enumerate}[label=\Alph*., itemsep=0pt]
    \item $\overrightarrow{AC}=\vec{a}+\vec{b}$
    \item $\overrightarrow{AE}=\vec{a}+\dfrac{1}{2}\vec{b}$
    \item $\overrightarrow{EF}=-\dfrac{1}{2}\vec{a}+\dfrac{1}{2}\vec{b}$
    \item $\overrightarrow{AE}\cdot\overrightarrow{AF}=\dfrac{3}{2}|\vec{a}|^2$
  \end{enumerate}

\item 已知 $|\vec{a}|=1$，$|\vec{b}|=2$，$\vec{a}$ 与 $\vec{b}$ 的夹角为 $\dfrac{\pi}{3}$，则（\quad）
  \begin{enumerate}[label=\Alph*., itemsep=0pt]
    \item $\vec{a}\cdot\vec{b}=1$
    \item $|\vec{a}+\vec{b}|=\sqrt{7}$
    \item $|\vec{a}-\vec{b}|=\sqrt{3}$
    \item $(\vec{a}+\vec{b})\perp(\vec{a}-\vec{b})$
  \end{enumerate}
\end{enumerate}

\vspace{6pt}

% ============ 三、填空题 ============
\noindent\rule{\textwidth}{0.4pt}\\[2pt]
{\large\bfseries 三、填空题}（本大题共4小题，每小题5分，共20分。）

\vspace{6pt}

\begin{enumerate}[resume]
\item 已知向量 $\vec{a}=(3,4)$，$\vec{b}=(-1,2)$，则 $\vec{a}\cdot\vec{b}=$ \underline{\hspace{2cm}}。

\item 已知 $\vec{a}=(2,-1)$，$\vec{b}=(-4,k)$，若 $\vec{a}\parallel\vec{b}$，则 $k=$ \underline{\hspace{2cm}}。

\item 已知向量 $\vec{a}=(3,0)$，$\vec{b}=(1,\sqrt{3})$，则 $\vec{a}$ 与 $\vec{b}$ 的夹角为 \underline{\hspace{2cm}}。

\item 已知 $|\vec{a}|=2$，$|\vec{b}|=3$，$\vec{a}$ 与 $\vec{b}$ 的夹角为 $120°$，则 $|2\vec{a}+\vec{b}|=$ \underline{\hspace{2cm}}。
\end{enumerate}

\vspace{6pt}

% ============ 四、解答题 ============
\noindent\rule{\textwidth}{0.4pt}\\[2pt]
{\large\bfseries 四、解答题}（本大题共6小题，共70分。解答应写出文字说明、证明过程或演算步骤。）

\vspace{6pt}

\begin{enumerate}[resume]
\item[17.]（本小题10分）\\
已知 $\vec{a}=(1,-2)$，$\vec{b}=(3,4)$。
\begin{enumerate}[(1)]
  \item 求 $2\vec{a}+\vec{b}$ 的坐标；
  \item 求 $\vec{a}\cdot\vec{b}$ 及 $\vec{a}$ 与 $\vec{b}$ 的夹角的余弦值。
\end{enumerate}

\vspace{6cm}

\item[18.]（本小题12分）\\
在 $\triangle ABC$ 中，$\overrightarrow{AB}=\vec{a}$，$\overrightarrow{AC}=\vec{b}$，$|\vec{a}|=4$，$|\vec{b}|=3$，$\angle BAC=60°$，$M$ 为 $BC$ 的中点。
\begin{enumerate}[(1)]
  \item 求 $\vec{a}\cdot\vec{b}$；
  \item 求 $|\overrightarrow{AM}|$；
  \item 求 $\overrightarrow{AB}\cdot\overrightarrow{AM}$。
\end{enumerate}

\vspace{6cm}

\item[19.]（本小题12分）\\
已知非零向量 $\vec{a}$、$\vec{b}$ 满足 $|\vec{a}|=2$，$\vec{a}$ 与 $\vec{b}$ 的夹角为 $\dfrac{\pi}{3}$，且 $|\vec{a}+\vec{b}|=\sqrt{7}$。
\begin{enumerate}[(1)]
  \item 求 $|\vec{b}|$ 的值；
  \item 求 $|2\vec{a}-\vec{b}|$ 的值。
\end{enumerate}

\vspace{6cm}

\item[20.]（本小题12分）\\
已知 $|\vec{a}|=2$，$|\vec{b}|=1$，$\vec{a}$ 与 $\vec{b}$ 的夹角为 $\dfrac{\pi}{3}$。
\begin{enumerate}[(1)]
  \item 求 $|\vec{a}+\vec{b}|$；
  \item 求向量 $\vec{a}+\vec{b}$ 与 $\vec{a}$ 的夹角。
\end{enumerate}

\vspace{6cm}

\item[21.]（本小题12分）\\
在平行四边形 $ABCD$ 中，$|\overrightarrow{AB}|=6$，$|\overrightarrow{AD}|=4$，$\angle BAD=60°$，$E$ 为 $BC$ 的中点，$M$ 在 $AC$ 上且 $\overrightarrow{AM}=\dfrac{1}{3}\overrightarrow{AC}$。
\begin{enumerate}[(1)]
  \item 用 $\overrightarrow{AB}$ 和 $\overrightarrow{AD}$ 表示 $\overrightarrow{AE}$ 和 $\overrightarrow{ME}$；
  \item 求 $\overrightarrow{AM}\cdot\overrightarrow{ME}$ 的值。
\end{enumerate}

\vspace{6cm}

\item[22.]（本小题12分）\\
已知非零向量 $\vec{a}$、$\vec{b}$ 满足 $(\vec{a}-2\vec{b})\perp\vec{a}$，$(\vec{a}+\vec{b})\perp\vec{b}$。
\begin{enumerate}[(1)]
  \item 求 $\vec{a}$ 与 $\vec{b}$ 的夹角；
  \item 若 $|\vec{a}|=2$，求 $|\vec{a}+\vec{b}|$。
\end{enumerate}

\vspace{6cm}
\end{enumerate}

\vfill
\begin{center}
\small\textcolor{gray}{本试卷仅供教学参考}
\end{center}

\end{document}
